% Cayley's Collected Mathematical Papers, Volume 10, Pages 51--100
% Transcribed from: collectedmathema10cayluoft.pdf
% Output: cayley_vol10_pages_051_100.tex

\documentclass[11pt,a4paper]{article}
\usepackage[utf8]{inputenc}
\usepackage[T1]{fontenc}
\usepackage[english]{babel}
\usepackage{amsmath,amssymb,amsthm}
\usepackage{geometry}
\usepackage{titlesec}
\usepackage{fancyhdr}
\usepackage{graphicx}
\usepackage{hyperref}

\geometry{margin=1in}

\theoremstyle{plain}
\newtheorem{theorem}{Theorem}
\newtheorem{lemma}[theorem]{Lemma}
\newtheorem{corollary}[theorem]{Corollary}

\theoremstyle{definition}
\newtheorem{definition}[theorem]{Definition}

\theoremstyle{remark}
\newtheorem{remark}{Remark}

\titleformat{\section}{\large\bfseries}{\thesection.}{0.5em}{}
\titleformat{\subsection}{\normalsize\bfseries}{\thesubsection.}{0.5em}{}

\pagestyle{fancy}
\fancyhf{}
\fancyhead[C]{\small Cayley's Collected Mathematical Papers --- Vol.~10, pp.~51--100}
\fancyfoot[C]{\thepage}
\renewcommand{\headrulewidth}{0.4pt}

\setlength{\parindent}{1em}
\setlength{\parskip}{0.3em}

\newcommand{\paper}[5]{%
\newpage
\begin{center}
{\large\bfseries #1}\\[0.5em]
{\small [#2]}\\[0.3em]
{\small [#3, vol.~#4 (#5)]}
\end{center}
\vspace{0.5em}
}

\begin{document}

\begin{center}
{\LARGE\bfseries Collected Mathematical Papers}\\[1em]
{\Large Arthur Cayley}\\[1em]
{\large Volume 10, Pages 51--100}
\end{center}

\vspace{2em}

\tableofcontents

% ============================================================
% PAPER 641 (Conclusion)
% ============================================================
\newpage
\section*{641. On the Flexure of a Spherical Surface.}
\addcontentsline{toc}{section}{641. On the Flexure of a Spherical Surface.}
\markboth{CAYLEY}{On the Flexure of a Spherical Surface}

[From the \textit{Proceedings of the London Mathematical Society}, vol.~viii. (1877),
pp.~12, 13.]

\textit{(Conclusion from p.~30.)}

\begin{center}
[Page~31]
\end{center}

of the quadrilateral $ABCD$ consider the birectangular triangle $ABE$;
the new form of this is $A'B'E'$, where the radius $OA' = k \cdot OA = k$,
is less than the original radius unity, but $OE' = k/c$, is greater than
the same radius unity.

The surface has at $E'$ a conical point, the semi-aperture of the cone
being $= \tan^{-1} \rho$, if as usual $k' = \sqrt{1-k^2}$;
to verify this, writing for convenience $q=0$, we have for the meridian
section $x = \cos\phi$, $z = E(k,\phi)$, and thence
\[
\frac{dz}{dx} = -\frac{\sqrt{1-k^2\sin^2\phi}}{\sin\phi},
\]
which for $\phi = 90^\circ$ gives $\dfrac{dz}{dx} = -k'$.
Observe also that for $\phi = 0$, $\dfrac{dz}{dx} = \infty$,
viz. the surface of revolution cuts the plane of the equator at right angles.

There is no limit to the arc $AB$, it may be $= 360^\circ$,
viz. we must in this case cut the hemisphere along a meridian to allow of
the deformation; or it may exceed $360^\circ$, the spherical surface being in
this case conceived of as wrapping indefinitely over itself, and we may
instead of the half lune $E'A'B'$, consider the lune included between
two meridians extending from pole to pole, and therefore the whole
spherical surface, conceived of as wrapping indefinitely over itself;
the result is, that this may be deformed into a surface of revolution,
which, in its general form, resembles that obtained by the revolution of
an arc less than a semi-circle round its chord; the half-chord being
greater, and the versed-sine less than the radius of the original sphere.

If $k > 1$, there is obviously a limit to the latitude
$AD = BC$, of the spherical quadrilateral; viz. this is equal to
$\sin^{-1}\tfrac{1}{k}$.
Supposing that in the quadrilateral $ABCD$ (fig.~1) the latitude has
this limiting value, then (see fig.~2) the new form is $A''B''C''D''$,
where along the bounding arc $C B''$ the tangent plane is horizontal;
viz. as before
\[
\frac{dz}{dx} = -\frac{\sqrt{1-k^2\sin^2\phi}}{\sin\phi} = 0
\quad \text{for } \phi = \sin^{-1}\frac{1}{k}.
\]
It is to be observed, that the radii for the parallels $A''B''$ and
$C D''$ are $k$ and $\tfrac{k}{c}\cos\phi$ respectively; the difference of

\begin{center}
[Page~32]
\end{center}

these is $k(1-\cos\phi)$, which, however great $k$ is, must be less than
the arc of meridian
\[
A''D'' = \phi;
\]
substituting for $k$ the value $\dfrac{1}{\sin\phi}$, the condition is
$\dfrac{\phi}{\sin\phi} > \phi$, viz. this is
\[
\tan\tfrac{1}{2}\phi < \phi,
\]
which is true for every value up to $\phi = 90^\circ$.

But, more than this, we should have
\[
k'(1-\cos\phi)^2 + E^2(k,\phi) < \phi^2,
\]
viz. writing as before $k = \dfrac{1}{\sin\phi}$, this is
\[
E^2(\phi', \phi) < \phi^2 - \frac{(1-\cos\phi)^2}{\sin^2\phi};
\]
this must be true, although (relating as it does to a form of $E$ for
which $k$ is greater than 1) there might be some difficulty in verifying it.

There is, as in the first case, no limit to the value of $AB$,
viz. this may be $= 360^\circ$, the spherical zone being then cut along a
meridian, or it may be greater than $360^\circ$; and, moreover, the spherical
quadrilateral may extend south of the equator, but of course so that the
limiting south latitude does not extend beyond the foregoing value
$\sin^{-1}\tfrac{1}{k}$:
viz. we may have a zone between the latitudes
$\pm \sin^{-1}\tfrac{1}{k}$, which may be a complete zone from longitude
$0^\circ$ to $360^\circ$ or to any greater value than $360^\circ$.

The result is, that the zone is deformed into a surface of revolution,
which in its general form resembles that obtained by the revolution of a
half-circle or half-ellipse about a line parallel to and beyond its
bounding diameter, the bounding half-diameter being less, and the
greatest radius of rotation greater, than the radius of the original sphere.


% ============================================================
% PAPER 642
% ============================================================
\newpage
\section*{642. On a Differential Relation Between the Sides of a Quadrangle.}
\addcontentsline{toc}{section}{642. On a Differential Relation Between the Sides of a Quadrangle.}
\markboth{CAYLEY}{On a Differential Relation Between the Sides of a Quadrangle}

[From the \textit{Messenger of Mathematics}, vol.~vi. (1877), pp.~99--101.]

\begin{center}
[Page~33]
\end{center}

Let the sides and diagonals $YZ$, $ZX$, $XY$, $OX$, $OY$, $OZ$ of a
quadrangle be $f$, $g$, $h$, $a$, $b$, $c$, and let the component triangles
be denoted as follows:
\[
A = \triangle YZO = (b,c,f), \quad
B = \triangle ZX O = (c,a,g), \quad
C = \triangle XY O = (a,b,h), \quad
\Omega = \triangle XYZ = (f,g,h),
\]
viz. $A$, $B$, $C$, $\Omega$ are the triangles whose sides are
$(b,c,f)$, $(c,a,g)$, $(a,b,h)$, $(f,g,h)$ respectively,
so that $\Omega = A + B + C$.

Then we have between $(a,b,c,f,g,h)$ an equation giving rise to a
differential relation, which may be written
\[
\Omega(Aa\,da + Bb\,db + Cc\,dc) - (BCf\,df + CAg\,dg + ABh\,dh) = 0.
\]
This may be proved geometrically and analytically.

First, for the geometrical proof, it is enough to prove that, when
$a$ and $b$ alone vary, the relation between the increments is
$Aa\,da + Bb\,db = 0$;
for then $a$ and $g$ alone varying, the relation between the increments
will be $Ca\,da - Cg\,dg = 0$
(as to the negative sign it is clear from the figure that $a$, $g$
will increase or diminish together):
and we thence at once infer the general relation.

We have consequently to prove that, considering $a$ and $b$ as alone
variable,
\[
Aa\,da + Bb\,db = 0;
\]
or, what is the same thing,
\[
Aa\,da : -Bb\,db = \triangle XOZ : \triangle YOZ.
\]
The points $XYZ$ remain fixed; but $O$ moves through the infinitesimal
arc $OO'$, centre $Z$, which may be considered as situate in the right
line $OM$ drawn from $O$ at right angles to $ZO$, and meeting $XY$
produced in the point $M$.

And then, writing for a moment $\angle OXY = X$, $\angle OYX = Y$,
$\angle OMY = M$, we find at once
\[
da = OO' \cos(X+M), \qquad -db = OO' \cos(Y-M);
\]
that is,
\[
\frac{da}{-db} = \frac{\cos(X+M)}{\cos(Y-M)},
\]
or
\[
\frac{Aa\,da}{Bb\,db} = \frac{a\cos(X+M)}{b\cos(Y-M)}.
\]
But drawing $X\alpha$, $Y\beta$ each of them at right angles to $ZO$, we
have $a\cos(X+M) = X\alpha$, $b\cos(Y-M) = Y\beta$, and evidently
$\triangle XOZ : \triangle YOZ = X\alpha : Y\beta$;
whence the equation is
\[
\frac{Aa\,da}{Bb\,db} = \frac{\triangle XOZ}{\triangle YOZ},
\]
which is the required relation.

\begin{center}
[Page~34]
\end{center}

For the analytical proof, it is to be observed that the relation
between $a$, $b$, $c$, $f$, $g$, $h$ is a quadric relation in the
quantities $a^2$, $b^2$, $c^2$, $f^2$, $g^2$, $h^2$ respectively;
this may be written
\begin{multline*}
1 + a^2(b^2 + c^2 - f^2)(g^2 + h^2 - f^2) \\
+ b^2(c^2 + a^2 - g^2)(h^2 + f^2 - g^2) \\
+ c^2(a^2 + b^2 - h^2)(f^2 + g^2 - h^2) \\
- (b^2 + c^2 - f^2)(c^2 + a^2 - g^2)(a^2 + b^2 - h^2) = 0;
\end{multline*}
say for a moment this is $A + Ba^2 + Ca^4 = 0$, where
\begin{align*}
A &= b^2g^2 + b^2h^2 + c^2h^2 + c^2g^2 + f^2(b^2 - h^2)(c^2 - g^2) \\
  &\quad - (b^2 + c^2)g^2h^2 - (g^2 + h^2)b^2c^2, \\
B &= -(b^2 - c^2)(g^2 - h^2) + f^2(-b^2 - c^2 + f^2 - g^2 - h^2) + f^4, \\
C &= c^2;
\end{align*}
then we have as usual
\[
\frac{du}{da}\,da + \frac{du}{db}\,db + \&\text{c.} = 0,
\]

\begin{center}
[Page~35]
\end{center}

where
\[
\frac{du}{da} = Ca^2 + B.
\]
But in virtue of $u = 0$, we have
\[
(Ca^2 + B)^2 = C(Ca^4 + Ba^2 + A) + (B^2 - AC),
\]
that is, $\dfrac{du}{da} = \sqrt{B^2 - 4AC}$;
and here $B^2 - 4AC$ is a quartic function of $f^2$, which is easily seen
to reduce itself to the form
\[
f^2 - (g+h)^2, \quad f^2 - (g-h)^2, \quad f^2 - (b+c)^2, \quad f^2 - (b-c)^2.
\]
The coefficients of $da$, $db$, $dc$, \&c., are given as expressions of
the like form; substituting their values, the differential relation is
\[
\sqrt{[f^2 - (g+h)^2][f^2 - (g-h)^2][f^2 - (b+c)^2][f^2 - (b-c)^2]}\, Aa\,da
+ \&\text{c.} = 0,
\]
which is, in fact, the foregoing result.

It is right to notice that there are in all 16 linear factors,
\begin{alignat*}{3}
&f+g+h, \quad && b+c+f, \quad && c+a+g, \quad a+b+h \\
& \text{say } d, && f, && g, \quad h, \\
&-f+g+h, \quad && -b+c+f, \quad && -c+a+g, \quad -a+b+h \\
& \text{say } d', && f', && g', \quad h', \\
&\phantom{-}f-g+h, \quad && \phantom{-}b-c+f, \quad && \phantom{-}c-a+g, \quad \phantom{-}a-b+h \\
& \text{say } d'', && f'', && g'', \quad h'', \\
&\phantom{-}f+g-h, \quad && \phantom{-}b+c-f, \quad && \phantom{-}c+a-g, \quad \phantom{-}a+b-h \\
& \text{say } d''', && f''', && g''', \quad h''',
\end{alignat*}
and this being so, the coefficients of $Aa\,da$, $Bb\,db$,
$Cc\,dc$, $Ff\,df$, $Gg\,dg$, $Hh\,dh$, are
\begin{align*}
&\sqrt{dd'd''d''' \cdot ff'f''f''' \cdot gg'g''g''' \cdot hh'h''h'''}, \\
&\sqrt{dd'd''d''' \cdot gg'g''g''' \cdot hh'h''h''' \cdot ff'f''f'''}, \\
-&\sqrt{hh'h''h''' \cdot ff'f''f''' \cdot gg'g''g''' \cdot dd'd''d'''}, \&\text{c.}
\end{align*}
respectively.

We may imagine the quadrilateral $ZOXY$ composed of the four rods
$ZO$, $OX$, $XY$, $YZ$ (lengths $c$, $a$, $h$, $f$ as before) jointed
together at the angles, and kept in equilibrium by forces $B$, $G$
acting along the diagonals $OY (=b)$, $ZX (=g)$ respectively.
We have $c$, $a$, $h$, $f$ given constants, and the relation
$\phi(a,b,c,f,g,h) = 0$, which connects the six quantities is the
relation between the two variable diagonals $(g,b)$; by what precedes,
the differential relation $\phi'_g \cdot dg + \phi'_b \cdot db = 0$
is equivalent to $\Omega Bb\,db - CAg\,dg = 0$.

By virtual velocities we have as the condition of equilibrium
$B\,db + G\,dg = 0$;
hence, eliminating $db$, $dg$ we have
\[
\frac{B}{\Omega Bb} = \frac{G}{-CAg},
\quad\text{or, say}\quad
\frac{B}{b} = \frac{G}{g} = \frac{1}{\triangle XYO \cdot \triangle ZYO
- \triangle XYZ \cdot \triangle ZXO},
\]
viz. the forces, divided by the diagonals along which they act, are
proportional to the reciprocals of the products of the two pairs of
triangles which stand on these diagonals respectively.
The negative sign shows, what is obvious, that the forces must be, one
of them a pull, the other a push.


% ============================================================
% PAPER 643
% ============================================================
\newpage
\section*{643. On a Quartic Curve with Two Odd Branches.}
\addcontentsline{toc}{section}{643. On a Quartic Curve with Two Odd Branches.}
\markboth{CAYLEY}{On a Quartic Curve with Two Odd Branches}

[From the \textit{Messenger of Mathematics}, vol.~vi. (1877), pp.~107, 108.]

\begin{center}
[Page~36]
\end{center}

It is a known theorem that the branches of a plane curve are even or odd;
viz. two even branches, or an even and an odd branch (whether of the
same curve or of different curves) intersect in an even number (it may
be 0, and this is to be understood throughout) of real points; but two
odd branches (of the same curve or of different curves) intersect in an
odd number of real points\footnote{The two branches must be distinct branches; a branch whether odd or
even does not of necessity intersect itself (have upon it any real node),
but it may intersect itself in an odd, or an even, number of real points.}.

In particular, a right line is an odd branch, and hence it meets any
even branch of a curve in an even number of real points, and an odd
branch in an odd number of real points; or (what is the same thing) an
even branch is one which is met by any right line whatever in an even
number of real points; and an odd branch is one that is met by any
right line whatever in an odd number of real points.

It is to be observed, that the simple term branch is used to denote
what has been called a complete branch, viz. the partial branches which
touch an asymptote at its opposite extremities are considered as parts
of one and the same branch, and so in other cases. Thus a quadric
curve, whether ellipse, parabola, or hyperbola, is one even branch; a
cubic curve is either one odd branch, or else it is an odd branch and
even branch; and generally a curve of an odd order has always an odd
number of odd branches, and a curve of an even order has always an
even number of odd branches.

A curve without nodes has at most one odd branch; for if there were
two, these would intersect in a real point, which would be a real node
on the curve. In particular, a quartic curve having two odd branches
must have a real node; this however may be, as in the instance about
to be given, a node at infinity.

A simple instance of a quartic curve with two odd branches is that
represented by the equation
\[
(a^2 - 1)(y^2 + 1) - 2mxy = 0,
\]

\begin{center}
[Page~37]
\end{center}

or, what is the same thing,
\[
x^2 = \tfrac{1}{2}\{2 + m^2 + m\sqrt{4 + m^2}\}, \qquad
\frac{1}{a^2} = \tfrac{1}{2}\{2 + m^2 - m\sqrt{4 + m^2}\},
\]
so that $m$ being positive $a > 1$, and the curve consists of two real
branches included between the lines $x = a$, $x = \tfrac{1}{a}$, and the
lines $x = -a$, $x = -\tfrac{1}{a}$ respectively; each of these lines
touches the curve in a real point, viz. $x$ having any one of the
last-mentioned values, the value of $y$ at the point of contact is
$y = \tfrac{m}{2x}$; and between each pair of lines we have the asymptote
$x = +1$ or $x = -1$.

Hence the curve has the form shown in the figure, and it is thereby
evident, that each branch of the curve is met by any real right line
whatever in one real point, or else in three real points.

The numerical values in the figure are $a = \tfrac{3}{2}$, $m = \tfrac{5}{4}$,
whence also $x = a$ or $-\tfrac{1}{a}$, $y = 1$, and $x = -a$ or
$\tfrac{1}{a}$, $y = -1$.

The curve has two nodes at infinity, viz. writing the equation in the form
\[
(x^2 - z^2)(y^2 + z^2) - 2mxyz^2 = 0,
\]
that is,
\[
x^2y^2 + z^2(x^2 - y^2 - 2mxy) + z^4 = 0,
\]
it appears that the points $(z=0, x=0)$, $(z=0, y=0)$ are each of them a node.
The first of these $(z=0, x=0)$ is the real intersection of the two odd
branches: the other of them is a conjugate point.


% ============================================================
% PAPER 644
% ============================================================
\newpage
\section*{644. Note on Magic Squares.}
\addcontentsline{toc}{section}{644. Note on Magic Squares.}
\markboth{CAYLEY}{Note on Magic Squares}

[From the \textit{Messenger of Mathematics}, vol.~vi. (1877), p.~168.]

\begin{center}
[Page~38]
\end{center}

In a magic square of any odd order, formed according to the ordinary
process, there is a tolerably simple analytical expression for the number
which occupies any given compartment; thus taking the square of 21, let
the dexter diagonals (N.W. to S.E.) commencing from the N.E. corner
compartment, be numbered $1$, $2$, $3$, \dots, $20$, $21$, $20'$, $19'$,
\dots, $2'$, $1'$, the diagonals of course containing these numbers of
compartments respectively; and in any diagonal let the compartments
reckoning from the top line be numbered $1$, $2$, $3$, \dots\ respectively;
then if $D_{\theta,\phi}$ (or $D'_{\theta,\phi}$ as the case may be) denotes
the number in the compartment $\phi$ of the diagonal $\theta$ or $\theta'$,
we have
\begin{align*}
D_{\theta+\phi,\phi} &= 20(\theta+\phi) + 10 + \phi, \\
D_{\theta',\phi} &= 20\theta' + 231 + \phi(-21), \\
D_{\theta'+\phi',\phi} &= -22\phi + 430 + \phi,
\end{align*}
where $\theta + \phi$ is to be replaced by its value, $= \theta'$ or $> \theta'$ as the
case may be; and in like manner for the other half of the square; it
would not be difficult, but it would take some space, to give the general
formulae for a square of any odd order.

The formula for a square of the order $2n+1$ is at once obtained from
the consideration that the numbers of the several dexter diagonals
commence with the terms of the arithmetical progression
\[
2n^2 + n, \quad 2n^2 + 3n + 1, \quad 2n^2 + 5n + 2, \quad \dots,
\]
viz. the dexter diagonal through the N.E. corner compartment (which
contains a single compartment, No. $2n^2 + n + 1$) is called diagonal $1$;
the next dexter diagonal (which contains two compartments numbered
$2n^2 + n$, $2n^2 + n + 2$) is called diagonal $2$; and so on.


% ============================================================
% PAPER 645
% ============================================================
\newpage
\section*{645. A Smith's Prize Paper, 1877.}
\addcontentsline{toc}{section}{645. A Smith's Prize Paper, 1877.}
\markboth{CAYLEY}{A Smith's Prize Paper, 1877}

[From the \textit{Cambridge Examination Papers}, 1877.]

\begin{center}
[Page~39]
\end{center}

\begin{enumerate}
\item Given a curve of the order $m$, and a line meeting it in $m$ real
points $O_1$, $O_2$, \dots, $O_m$, find the number of the tangents from
the several points $O_1$, $O_2$, \dots, $O_m$, which lie altogether in the
plane of the curve; and thence infer the theorem ``If a line meet a
curve of the order $m$ in $m$ real points, the number of real tangents
from these points to the curve is $m(m-1)$ or $m(m-2)$ according as $m$
is odd or even.''

\item Show that the locus of a point $P$, such that its polars with
respect to three conics meet in a point, is a curve of the third order.

\item Find the condition which must be satisfied in order that the equation
\[
(a,b,c,d\between x,y)^3 = 0
\]
may represent three lines; and in particular when the three lines are
concurrent. Find also the condition that the equation may represent a
line and its two tangentials (viz. lines each of them meeting the line
in an inflexion of the cubic curve represented by the equation).

\item In the cubic curve
\[
y^2 = 4x^3 - g_2x - g_3,
\]
if $x$, $y$ be expressed as functions of the parameter $u$, show that
\[
x(u+v) + x(u) + x(v) = \frac{1}{4}\left(
\frac{y(u) - y(v)}{x(u) - x(v)} \right)^2.
\]

\item Find an expression for the number of the conics which pass through
four given points and touch a given line; and account for the different
kinds of solution.

\item Find the differential equation of the lines of curvature of an
ellipsoid; and (without attempting a general solution) show that through
a given point on the ellipsoid there pass four umbilici.
\end{enumerate}

\begin{center}
[Page~40]
\end{center}

\begin{enumerate}
\setcounter{enumi}{6}
\item Find the equation of the surface which is the locus of the
foot of the perpendicular from the centre of an ellipsoid on the tangent
plane at any point of one of its lines of curvature; and show that this
surface is a particular case of the wave surface.

\item An ellipsoid and a hyperboloid of one sheet are confocal; taking
the equations to be
\[
\frac{x^2}{a^2} + \frac{y^2}{b^2} + \frac{z^2}{c^2} = 1, \qquad
\frac{x^2}{a^2+\theta} + \frac{y^2}{b^2+\theta} + \frac{z^2}{c^2+\theta} = 1,
\]
show that the intersection consists of curves of curvature on each
surface respectively; and give a construction (depending on the
rectilinear generation of the hyperboloid) for finding at any point of
the intersection the normal of the ellipsoid.

\item The consecutive positions of a rigid body are determined by the
values of six parameters; assuming the general expressions for the
sliding and rotating velocities, form the condition that the motion may
be simply rotational; and explain the result.

\item Find the condition in order that the equation
\[
a\cos x + b\cos y + c\cos z + d\cos w + e\cos t + f = 0
\]
may be satisfied identically (i.e.\ for all values of $x$, $y$, $z$, $w$, $t$);
and thence find the condition in order that the equation
\[
a\cos x + b\cos y + c\cos z + d\cos w + e\cos t + f = 0
\]
may have a simultaneous solution with its derived equation
\[
a\sin x + b\sin y + c\sin z + d\sin w + e\sin t = 0.
\]
\end{enumerate}

\begin{center}
[Page~41]
\end{center}

\begin{enumerate}
\setcounter{enumi}{10}
\item Find the orthogonal trajectory of the series of parabolas
\[
y^2 = 4a(x+a),
\]
where $a$ is the variable parameter.

\item Explain the different singularities of a curve; and in particular
show that a cusp is to be regarded as a double point.

\item Show how to find the equation of the first or second reciprocal
of a given curve; and apply it to find the reciprocal of a cubic with
respect to a conic.
\end{enumerate}

\begin{center}
[Page~42]
\end{center}

\begin{enumerate}
\setcounter{enumi}{13}
\item Find the number of the lines which meet each of four given lines
and touch a given quadric surface.

\item Two surfaces of the orders $m$, $n$ respectively intersect in a
curve of the order $\mu$; find the order of the cone standing on this
curve, the vertex of the cone being a given point; and thence find the
order of the scroll generated by the lines which meet the curve of
intersection in two points, and also meet a given line.

\item Find the partial differential equation of the surfaces which cut
at right angles a given series of surfaces
\[
F(x,y,z) = C.
\]
\end{enumerate}

\begin{center}
[Page~43]
\end{center}

\begin{enumerate}
\setcounter{enumi}{16}
\item In the question of the attraction of an ellipsoid on an external
point, assuming only the attraction of a parallel plate on an external
point, and that of a solid elliptic disc on a point in its axis, show
how the result is obtained.

\item Find the attraction of a solid ellipsoid on an external point
situate on one of the principal sections, the principal axes being all
unequal.

\item A body moves under the action of a central force; find the
differential equation of the orbit; and show that the orbit is a conic
section when the force varies as the inverse square of the distance.
\end{enumerate}

\begin{center}
[Page~44]
\end{center}

\begin{enumerate}
\setcounter{enumi}{19}
\item The roots of the equation
\[
\frac{x^2}{a^2+\theta} + \frac{y^2}{b^2+\theta} + \frac{z^2}{c^2+\theta} = 1
\]
are $\lambda$, $\mu$, $\nu$; explain the geometrical signification of
$(\lambda,\mu,\nu)$ and $(\lambda,\mu)$ respectively; and find the
partial differential equation of the surfaces which cut at right angles
the series of confocal quadrics
\[
\frac{x^2}{a^2+\lambda} + \frac{y^2}{b^2+\lambda} + \frac{z^2}{c^2+\lambda} = 1.
\]

\item Find the potential of the attraction of a solid ellipsoid on an
external point; and thence the component attractions.

\item Find the attraction of a solid elliptic disc on an external
point in its axis; and show that it is the same as if the mass were
condensed into a focal conic.
\end{enumerate}

\begin{center}
[Page~45]
\end{center}

\begin{enumerate}
\setcounter{enumi}{22}
\item Find the differential equation of the asymptotic lines on a
given surface; and apply it to determine these lines on the helicoid
and on the catenoid.

\item Show that the two sets of lines of curvature are at right angles
to each other; and find the differential equation of the lines of
curvature on the paraboloid
\[
\frac{x^2}{a} + \frac{y^2}{b} = 2z.
\]

\item Find the equation of the surface of the wave; and show that it
is the same surface as that which is the locus of the feet of the
perpendiculars from the centre on the tangent planes to the lines of
curvature of an ellipsoid.
\end{enumerate}

\begin{center}
[Page~46]
\end{center}

\begin{enumerate}
\setcounter{enumi}{25}
\item Find the general expression for the curvature of a plane curve;
and in particular for the curve $y = f(x)$.

\item Find the equation of a geodesic line on a surface; and show that
on a developable surface the geodesic lines become right lines when the
surface is developed into a plane.

\item Explain what is meant by the ``Index of a Point;'' and find the
index of each of the conical points of the wave surface.
\end{enumerate}


% ============================================================
% PAPER 646
% ============================================================
\newpage
\section*{646. On the General Equation of Differences of the Second Order.}
\addcontentsline{toc}{section}{646. On the General Equation of Differences of the Second Order.}
\markboth{CAYLEY}{On the General Equation of Differences of the Second Order}

[From the \textit{Quarterly Journal of Pure and Applied Mathematics},
vol.~xiv. (1877), pp.~79--85.]

\begin{center}
[Page~47]
\end{center}

I consider the general equation of differences of the second order,
\[
u \xi_{n+2} + B \xi_{n+1} + C \xi_n = 0,
\]
where $A$, $B$, $C$ are given constant coefficients. Writing this in
symbolical form as
\[
(A\varpi^2 + B\varpi + C)\xi_n = 0,
\]
where $\varpi$ denotes the operation of passing from $\xi_n$ to $\xi_{n+1}$,
the solution depends on the roots $\alpha$, $\beta$ of the quadratic
equation $A\varpi^2 + B\varpi + C = 0$.

The case of equal roots, $\alpha = \beta$, gives rise to the solution
\[
\xi_n = (A' + B'n)\alpha^n,
\]
and the general case where the roots are unequal gives
\[
\xi_n = A'\alpha^n + B'\beta^n.
\]
The paper goes on to consider the general theory, and in particular the
relation of the solution to the continued fraction
\[
\frac{C}{B+}\frac{AC}{B+}\frac{AC}{B+}\cdots
\]
and the corresponding series
\[
u \xi_0 - B\xi_1 + C\xi_2 - \dots.
\]

The complete theory requires the consideration of the cases where the
roots are equal, and where they are unequal; and the relation of the
solution to the particular integrals.

\begin{center}
[Page~48]
\end{center}

The analytical theory of the equation of differences is continued,
showing the relation between the successive coefficients and the
convergents of the continued fraction. If we write
\[
\frac{C}{B+}\frac{AC}{B+}\frac{AC}{B+}\cdots =
\frac{p_n}{q_n},
\]
then $p_n$ and $q_n$ satisfy the same equation of differences, with
different initial conditions. The general theory is established for
the solution of the equation in terms of the roots $\alpha$, $\beta$.

\begin{center}
[Page~49]
\end{center}

The solution is thus completed, and the relation to the theory of
continued fractions is fully exhibited. The different forms of the
solution are considered, according as the roots are real or complex,
and the particular case of equal roots is treated as a limiting case
of the general theory.


% ============================================================
% PAPER 647
% ============================================================
\newpage
\section*{647. On the Quartic Surfaces Represented by the Equation, Symmetrical Determinant $=0$.}
\addcontentsline{toc}{section}{647. On the Quartic Surfaces Represented by the Equation, Symmetrical Determinant $=0$.}
\markboth{CAYLEY}{On Quartic Surfaces Represented by the Equation, Symmetrical Determinant $=0$}

[From the \textit{Quarterly Journal of Pure and Applied Mathematics},
vol.~xiv. (1877), pp.~86--101.]

\begin{center}
[Page~50]
\end{center}

I consider the quartic surface represented by the equation
\[
\begin{vmatrix}
a & h & g & l \\
h & b & f & m \\
g & f & c & n \\
l & m & n & d
\end{vmatrix} = 0,
\]
where the coordinates are subjected to the condition that the
symmetrical determinant vanishes. This is a particular form of quartic
surface, having properties connected with the theory of quadrics and
the geometry of the position.

\begin{center}
[Page~51]
\end{center}

The surface is considered in relation to the theory of the quadric
locus represented by the equation
\[
(a,b,c,d,f,g,h,l,m,n\between x,y,z,w)^2 = 0,
\]
and the several forms which the quartic surface may assume, according
to the different relations between the coefficients.

\begin{center}
[Page~52]
\end{center}

The equation $V = 0$ thus represents a quartic surface; and the
different forms are considered according to the nature of the
component quadrics. The surface has in general nodes or other
singularities, which depend on the particular relations between the
coefficients of the fundamental quadric.

\begin{center}
[Page~53]
\end{center}

For each set of values of the coordinates satisfying the equation,
there corresponds a point on the surface; and the particular case where
the determinant reduces to the product of linear factors gives rise to
the consideration of the conic in question.

\begin{center}
[Page~54]
\end{center}

For each of the possible configurations, the surface may be regarded
as the envelope of a quadric, or as the locus of a singular point on
a system of quadrics. The theory is closely connected with that of
the bitangent planes and the nodal curves of the surface.

\begin{center}
[Page~55]
\end{center}

This is further developed; the several species of the surface are
classified according to the nature of the singularities. The
analytical conditions for the different forms are given, and the
geometrical interpretation is supplied.

\begin{center}
[Page~56]
\end{center}

The different species are tabulated, and the particular cases of
degeneracy are considered. The surface may reduce to a quadric counted
twice, or to the combination of two quadrics, or to other simpler
forms, according to the relations between the coefficients.


% ============================================================
% PAPER 648
% ============================================================
\newpage
\section*{648. Algebraical Theorem.}
\addcontentsline{toc}{section}{648. Algebraical Theorem.}
\markboth{CAYLEY}{Algebraical Theorem}

[From the \textit{Quarterly Journal of Pure and Applied Mathematics},
vol.~xiv. (1877), p.~102.]

\begin{center}
[Page~57]
\end{center}

If $a$, $b$, $c$, $d$, $e$, $f$, $g$, $h$, $i$, $j$, $k$, $l$, $m$, $n$, $o$, $p$
are the elements of a determinant of the 4th order, then the determinant
can be expressed in terms of its minors; and there exist relations
between the minors of different orders which are expressed by the
identities connecting the complementary minors. The theorem relates
to the expression of the square of any minor of order 2 as a function
of the minors of orders 1 and 3.

In particular, if we consider the minors formed by omitting any two
rows and columns, there exist relations of the form
\[
\begin{vmatrix}
A_1 & B_1 & C_1 & D_1 \\
A_2 & B_2 & C_2 & D_2 \\
A_3 & B_3 & C_3 & D_3 \\
A_4 & B_4 & C_4 & D_4
\end{vmatrix} = 0,
\]
where the $A$'s, $B$'s, \&c.\ denote minors of the original determinant.


% ============================================================
% PAPER 649
% ============================================================
\newpage
\section*{649. Addition to Mr Glaisher's Note on Sylvester's Paper, ``Development of an Idea of Eisenstein.''}
\addcontentsline{toc}{section}{649. Addition to Mr Glaisher's Note on Sylvester's Paper.}
\markboth{CAYLEY}{Addition to Mr Glaisher's Note on Sylvester's Paper}

[From the \textit{Quarterly Journal of Pure and Applied Mathematics},
vol.~xiv. (1877), pp.~103--105.]

\begin{center}
[Page~58]
\end{center}

The remark which I wish to make relates to the development of an idea
of Eisenstein, as presented in Sylvester's paper on the subject. The
point is that the formulae for the multiplication of the quotients of
the periods of the elliptic functions admit of a more general
presentation, in which the original forms are included as particular
cases.

If we consider the standard form of the relation between the
multipliers and the moduli, we have a system of equations which can be
written in the symmetrical form
\[
\sqrt{kk'} + \sqrt{ll'} = 1,
\]
where $k$, $k'$, $l$, $l'$ are the moduli and complementary moduli of
the elliptic functions related by the multiplication theorem.

The general theory leads to the consideration of the functions
\[
\frac{\vartheta_2(0|\tau)}{\vartheta_3(0|\tau)}, \quad
\frac{\vartheta_4(0|\tau)}{\vartheta_3(0|\tau)},
\]
and the modular equations which connect the different values of these
functions.

\begin{center}
[Page~59]
\end{center}

where the summation extends over all integer values of $m$ and $n$,
which satisfy the condition $an^2 + 2bmn + cm^2 = N$, and where $a$,
$b$, $c$ are coefficients connected with the quadratic form. The result
is expressed as a series involving the divisors of $N$, and can be
written in the form
\[
\sigma(N) = \sum_{d|N} d,
\]
where the summation extends over all divisors of $N$. The general
theory includes as particular cases the formulae of Eisenstein and the
subsequent developments by Sylvester.

The relation to the theory of binary quadratic forms is indicated, and
the connection with the class-number and the theory of genera.


% ============================================================
% PAPER 650
% ============================================================
\newpage
\section*{650. On a Quartic Surface with Twelve Nodes.}
\addcontentsline{toc}{section}{650. On a Quartic Surface with Twelve Nodes.}
\markboth{CAYLEY}{On a Quartic Surface with Twelve Nodes}

[From the \textit{Quarterly Journal of Pure and Applied Mathematics},
vol.~xiv. (1877), pp.~106--116.]

\begin{center}
[Page~60]
\end{center}

I consider the quartic surface represented by the equation
\[
p(YZ + XW) + q(ZX + YW) + r(XY + ZW) = 0,
\]
where $p + q + r = 0$. This surface has 12 nodes, which are situate
at the intersections of certain planes, and has properties analogous
to those of the cubic surface with 4 nodes.

The more general equation
\[
p(YZ - XW) + q(ZX - YW) + r(XY - ZW) = 0
\]
is also considered; and the relation between the two forms is
investigated. The surface with the plus signs has the 12 nodes, while
the surface with the minus signs has 4 conical points and other
singularities.

The 12 nodes are situate in threes on the 8 faces of the fundamental
tetrahedron; and the planes of the faces meet the surface in conics
which pass through the nodes.

\begin{center}
[Page~61]
\end{center}

The condition for a node is that the partial derivatives of the
function with respect to $X$, $Y$, $Z$, $W$ shall all vanish. This
gives the equations
\[
rY + qZ + pW = 0, \quad
rX + pZ + qW = 0, \quad
qX + pY + rW = 0, \quad
pX + qY + rZ = 0.
\]
Eliminating $X : Y : Z : W$, we obtain the condition that the
determinant
\[
\begin{vmatrix}
0 & r & q & p \\
r & 0 & p & q \\
q & p & 0 & r \\
p & q & r & 0
\end{vmatrix} = 0,
\]
which is satisfied in virtue of $p + q + r = 0$.

\begin{center}
[Page~62]
\end{center}

viz. the surface having the 12 nodes is the original surface where
\[
p(YZ + XW) + q(ZX + YW) + r(XY + ZW) = 0, \qquad p + q + r = 0.
\]

The surface may be regarded as a particular case of the general quartic
with 16 nodes; and the reduction from 16 to 12 is effected by the
special relations between the coefficients.


% ============================================================
% PAPER 651
% ============================================================
\newpage
\section*{651. On a Special Surface of Minimum Area.}
\addcontentsline{toc}{section}{651. On a Special Surface of Minimum Area.}
\markboth{CAYLEY}{On a Special Surface of Minimum Area}

[From the \textit{Quarterly Journal of Pure and Applied Mathematics},
vol.~xiv. (1877), pp.~117--126.]

\begin{center}
[Page~63]
\end{center}

The general equation of a minimal surface is known to be
\[
(1 + q^2)r - 2pqs + (1 + p^2)t = 0,
\]
where $p = \dfrac{dz}{dx}$, $q = \dfrac{dz}{dy}$, $r = \dfrac{d^2z}{dx^2}$,
$s = \dfrac{d^2z}{dx\,dy}$, $t = \dfrac{d^2z}{dy^2}$.

I consider the particular case where the surface is determined by the
condition that it contains a system of circles; and I find that there
is a special surface of minimum area which satisfies this condition.

Taking the equations of the minimal surface in Monge's form,
\[
x = X + X_1, \quad y = Y + Y_1, \quad z = Z + Z_1,
\]
where $X$, $Y$, $Z$ are functions of a parameter $\alpha$, and $X_1$,
$Y_1$, $Z_1$ are functions of a parameter $\beta$, and where
\[
X'^2 + Y'^2 + Z'^2 = 0, \qquad X_1'^2 + Y_1'^2 + Z_1'^2 = 0,
\]
we have the general solution; and the special case is obtained by
particular assumptions for the forms of these functions.

\begin{center}
[Page~64]
\end{center}

The conditions for the surface to contain circles are investigated.
If the surface contains a singly infinite system of circles, then the
circles must be lines of curvature on the surface; and the condition
for this gives a differential equation which must be satisfied by the
functions $X$, $Y$, $Z$ and $X_1$, $Y_1$, $Z_1$.

\begin{center}
[Page~65]
\end{center}

The solution of the equation leads to the consideration of elliptic
functions. We have
\[
X' = \frac{1}{2}\left(\frac{1}{\wp'\alpha}\right), \quad
Y' = \frac{1}{2}\left(\frac{\wp\alpha}{\wp'\alpha}\right), \quad
Z' = \frac{1}{2}\left(\frac{\wp^2\alpha}{\wp'\alpha}\right),
\]
where $\wp\alpha$ denotes the Weierstrassian elliptic function; and
similar expressions for $X_1'$, $Y_1'$, $Z_1'$ in terms of $\beta$.

\begin{center}
[Page~66]
\end{center}

The complete determination of the surface requires the integration of
these expressions; and we find
\begin{align*}
X &= \frac{1}{2}\int \frac{d\alpha}{\wp'\alpha}, \\
Y &= \frac{1}{2}\int \frac{\wp\alpha\,d\alpha}{\wp'\alpha}, \\
Z &= \frac{1}{2}\int \frac{\wp^2\alpha\,d\alpha}{\wp'\alpha}.
\end{align*}
These integrals can be expressed in terms of the elliptic integrals of
the first and second kinds; and the surface is thus completely
determined.

\begin{center}
[Page~67]
\end{center}

The surface is thus a particular minimal surface, distinguished from
the general minimal surface by the property of containing a system of
circles. It is a surface of the 10th order, and has certain
singularities which are investigated in the paper. The complete theory
is connected with that of the tetrahedroid and the wave surface.


% ============================================================
% PAPER 652
% ============================================================
\newpage
\section*{652. On a Sextic Torse.}
\addcontentsline{toc}{section}{652. On a Sextic Torse.}
\markboth{CAYLEY}{On a Sextic Torse}

[From the \textit{Quarterly Journal of Pure and Applied Mathematics},
vol.~xiv. (1877), pp.~127--132.]

\begin{center}
[Page~68]
\end{center}

A torse (developable surface) of the sixth order is considered,
having its generating lines arranged in a particular manner with
respect to the fundamental planes of reference. The equation of the
torse is obtained in the form
\[
\text{(sextic equation in } X, Y, Z, W\text{)} = 0,
\]
which represents a developable surface, that is, a surface generated by
the tangents to a curve.

The particular sextic torse has its edge of regression a curve of the
sixth order; and the generating lines of the torse are the tangents to
this curve. The equation may be written in the form
\[
(x^2 + y^2 + z^2)^3 - k(x^3 + y^3 + z^3)^2 = 0,
\]
or in a more general homogeneous form
\[
(X^2 + Y^2 + Z^2)^3W^2 - k(X^3 + Y^3 + Z^3)^2 = 0.
\]

The edge of regression is the curve of intersection of this surface
with another surface of the same order; and the nature of the
cuspidal edge is investigated.

\begin{center}
[Page~69]
\end{center}

The two half-branches of the cuspidal edge are considered, and the
nature of the point of regression is analyzed. The curve has a point
at which the two branches meet, and the tangent line at this point is
common to both branches.

\begin{center}
[Page~70]
\end{center}

We have
\[
r^2 = 1 + p^2 + q^2,
\]
and thence also, if $\tan\alpha = 2p$, $\pm 2\sqrt{r^2-1}$, that is
\[
\cos\alpha = \frac{1}{\sqrt{1+4p^2}}, \quad
\sin\alpha = \frac{2p}{\sqrt{1+4p^2}}.
\]

The analytical theory is developed, showing the relation between the
curvature of the edge of regression and the generating lines of the
torse. The point of closest approach of consecutive generators is
determined, and the equation of the tangent plane at any point of the
torse is found.

\begin{center}
[Page~71]
\end{center}

The sextic torse is completely determined by the condition that its
cuspidal edge is a curve of the sixth order with a particular
singularity. The classification of the different species of sextic
torses is indicated, and the relation to the general theory of
developable surfaces is explained.


% ============================================================
% PAPER 653
% ============================================================
\newpage
\section*{653. On a Torse Depending on the Elliptic Functions.}
\addcontentsline{toc}{section}{653. On a Torse Depending on the Elliptic Functions.}
\markboth{CAYLEY}{On a Torse Depending on the Elliptic Functions}

[From the \textit{Quarterly Journal of Pure and Applied Mathematics},
vol.~xiv. (1877), pp.~133--141.]

\begin{center}
[Page~73]
\end{center}

I consider a torse (developable surface) which depends on the elliptic
functions; viz. the equation of the torse involves the elliptic
functions of a parameter, and the cuspidal edge of the torse is a
transcendental curve.

The general theory of such torses requires the consideration of the
Weierstrassian elliptic functions $\wp u$, $\wp' u$, and their
connection with the geometry of the surface.

If we take the standard form
\[
y^2 = 4x^3 - g_2x - g_3,
\]
then the coordinates of a point on the cuspidal edge can be expressed as
\[
x = \wp u, \quad y = \wp' u, \quad z = u \wp u - 2\zeta u,
\]
where $\zeta u$ is the Weierstrassian zeta-function. These expressions
give the parametric representation of the edge of regression.

\begin{center}
[Page~74]
\end{center}

The conditions that the surface may be a torse (i.e.\ developable) are
that the consecutive generating lines shall intersect; and this gives
rise to a relation between the parameter and the coordinates which is
expressed by means of the elliptic functions.

The generating line at any point of the edge of regression is the
tangent to the curve; and the equations of this tangent can be written
in the form
\[
\frac{X - x}{\wp'' u} = \frac{Y - y}{\wp''' u} = \frac{Z - z}{1}.
\]

\begin{center}
[Page~75]
\end{center}

The analytical theory is developed, showing the relation of the torse
to the elliptic functions. The equation of the torse is obtained by
eliminating the parameter $u$ between the equations of the tangent
line and the condition of developability.

\begin{center}
[Page~76]
\end{center}

The equation of the torse is thus obtained in a form involving the
elliptic functions; and the singularities of the surface are
investigated. The cuspidal edge is a curve of the sixth order, and
the surface has also nodal and other singular lines.

\begin{center}
[Page~77]
\end{center}

We have thus the complete theory of the torse depending on the
elliptic functions; and the relation to the general theory of
developable surfaces is established. The particular cases where the
torse degenerates are considered, and the classification of the
different species is given.


% ============================================================
% PAPER 654 (starts on p. 79)
% ============================================================
\newpage
\section*{654. On Certain Octic Surfaces.}
\addcontentsline{toc}{section}{654. On Certain Octic Surfaces.}
\markboth{CAYLEY}{On Certain Octic Surfaces}

[From the \textit{Quarterly Journal of Pure and Applied Mathematics},
vol.~xiv. (1877), pp.~142--150.]

\begin{center}
[Page~79]
\end{center}

I consider certain surfaces of the eighth order (octic surfaces)
which arise in the theory of quadrics and the geometry of the
associated lines. The general equation of such a surface may be
written in the form
\[
\Phi(X,Y,Z,W) = 0,
\]
where $\Phi$ is a homogeneous function of the eighth degree.

If in this equation we write $c''h' = c$, $c'h'' = h$; the equation
takes a more particular form, and the surface has properties connected
with the configuration of the fundamental lines and points.

The octic surfaces considered are related to the theory of the
complex of lines and the congruences which are formed by the lines
meeting certain given curves. The surface is the locus of points
satisfying certain conditions with respect to the configuration.

\begin{center}
[Page~80]
\end{center}

The complete investigation of the surface requires the consideration
of the several species, according to the nature of the fundamental
configuration. The surface may have nodes and other singularities,
which correspond to particular positions of the generating lines.

The equation of the surface may be written in the form
\[
\begin{vmatrix}
A & H & G \\
H & B & F \\
G & F & C
\end{vmatrix} = 0,
\]
where $A$, $B$, $C$, $F$, $G$, $H$ are functions of the coordinates of
the second degree. This represents an octic surface which is the
envelope of a certain quadric locus.

The theory is continued in subsequent pages, with the investigation
of the particular forms and singularities of the surface.

\end{document}
